\documentclass[a4paper,12pt]{article}

\usepackage[T2A]{fontenc}
\usepackage[utf8]{inputenc}
\usepackage[russian]{babel}
\usepackage{amsmath,amssymb}
\usepackage{paratype}
\renewcommand{\familydefault}{\sfdefault}
\usepackage{icomma}
\usepackage[a4paper,margin=20mm,headheight=15pt]{geometry}
\usepackage{microtype}
\usepackage{tikz}
\usetikzlibrary{calc,arrows.meta,positioning,shapes.geometric}
\usepackage{tabularx,booktabs}
\usepackage{enumitem}
\usepackage{parskip}
\usepackage{fancyhdr}
\usepackage{setspace}
\usepackage{xcolor}

\definecolor{accent}{HTML}{008080}
\definecolor{darkaccent}{HTML}{004D4D}
\definecolor{textgray}{HTML}{333333}
\definecolor{softgray}{HTML}{6B7280}

\onehalfspacing
\setlength{\parindent}{0pt}
\setlength{\parskip}{0.5em}
\setlength{\emergencystretch}{2em}
\raggedbottom
\clubpenalty=10000
\widowpenalty=10000
\displaywidowpenalty=10000
\color{textgray}

\pagestyle{fancy}
\fancyhf{}
\fancyhead[L]{\small Ксения Васильченко · 6 класс}
\fancyhead[R]{\small НОД и НОК}
\fancyfoot[C]{\small\thepage}
\renewcommand{\headrulewidth}{0.4pt}
\renewcommand{\footrulewidth}{0pt}

\newcommand{\nodru}{\mathop{\text{НОД}}\nolimits}
\newcommand{\nokru}{\mathop{\text{НОК}}\nolimits}
\newcommand{\lessondate}{24 сентября 2026 г.}

\newcolumntype{Y}{>{\raggedright\arraybackslash}X}

\newcommand{\sectionline}{%
  \vspace{-0.2em}\par\noindent
  {\color{accent}\rule{\linewidth}{0.55pt}}%
  \vspace{0.35em}
}

\newcommand{\keyidea}[1]{%
  \par\noindent{\color{darkaccent}\bfseries #1}\par
}

\tikzset{
  flowstart/.style={
    ellipse,
    draw=accent,
    line width=0.9pt,
    align=center,
    minimum width=38mm,
    minimum height=10mm,
    font=\small
  },
  flowaction/.style={
    rounded corners=2mm,
    rectangle,
    draw=accent,
    line width=0.9pt,
    align=center,
    text width=46mm,
    minimum height=13mm,
    font=\small
  },
  flowdecision/.style={
    diamond,
    aspect=2.35,
    draw=darkaccent,
    line width=0.9pt,
    align=center,
    text width=43mm,
    inner sep=1.6pt,
    font=\small
  },
  flowresult/.style={
    rounded corners=2mm,
    rectangle,
    draw=darkaccent,
    line width=1pt,
    align=center,
    text width=39mm,
    minimum height=13mm,
    font=\small
  },
  flowarrow/.style={
    -{Stealth[length=3mm,width=2mm]},
    line width=0.9pt,
    draw=textgray
  }
}

\begin{document}

% ============================================================
% ТИТУЛЬНЫЙ ЛИСТ
% ============================================================

\thispagestyle{empty}
\begin{titlepage}
\centering
\vspace*{27mm}

{\Large\color{darkaccent}\bfseries Чек-лист\par}
\vspace{4mm}

{\huge\bfseries
НОД и НОК: вычисление и выбор в задачах\par}

\vspace{10mm}

{\large Ксения Васильченко, 6 класс\par}

\vfill

{\large\bfseries Что повторить после занятия\par}
\vspace{4mm}

\begin{minipage}{0.80\textwidth}
\centering
Десятичные и обыкновенные дроби · простейшие уравнения ·
разложение на простые множители · наибольший общий делитель ·
наименьшее общее кратное · выбор НОД или НОК по смыслу задачи
\end{minipage}

\vfill

{\normalsize
Лёвин Артём Александрович эксклюзивно для Ксении Васильченко\par}
\vspace{2mm}
{\normalsize \lessondate\par}

\vspace{18mm}

\begin{tikzpicture}[remember picture,overlay]
\draw[
  accent!55,
  line width=1.05pt
]
  ([xshift=20mm,yshift=17mm]current page.south west)
  .. controls
  ([xshift=67mm,yshift=30mm]current page.south west)
  and
  ([xshift=-64mm,yshift=8mm]current page.south east)
  ..
  ([xshift=-20mm,yshift=19mm]current page.south east);
\end{tikzpicture}

\end{titlepage}

% ============================================================
% ОСНОВНОЙ ЧЕК-ЛИСТ
% ============================================================

\section*{1. Короткая разминка}
\sectionline

Эти навыки использовались перед основной темой занятия.

\textbf{Десятичные дроби.}
При сложении и вычитании записывайте числа так, чтобы запятые стояли друг под другом:
\[
1{,}50+2{,}57=4{,}07,
\qquad
2{,}57-1{,}50=1{,}07.
\]
При умножении сначала умножайте числа без запятых, затем отделяйте справа столько цифр,
сколько знаков после запятой было в обоих множителях вместе:
\[
1{,}5\cdot1{,}2=1{,}80=1{,}8.
\]
При делении на натуральное число следите за положением запятой:
\[
1{,}5:3=0{,}5.
\]

\textbf{Обыкновенные дроби.}
Для сложения нужен общий знаменатель:
\[
\frac23+\frac17=\frac{14}{21}+\frac3{21}=\frac{17}{21}.
\]
При делении на дробь умножайте на обратную дробь:
\[
\frac23:3\frac15
=\frac23:\frac{16}{5}
=\frac23\cdot\frac5{16}
=\frac5{24}.
\]
Перед умножением дробей можно сократить общий множитель числителя одной дроби и знаменателя другой. Например, в произведении
\[
\frac23\cdot\frac5{16}
\]
можно сократить числа \(2\) и \(16\) на \(2\).

\textbf{Простейшие уравнения.}
\[
\frac{x}{3}=5 \quad\Rightarrow\quad x=15,
\]
а в уравнении
\[
\frac3x=5
\]
обязательно учитываем ограничение \(x\ne0\). Здесь
\[
x=\frac35.
\]
После нахождения корня выполните подстановку в исходное уравнение: левая и правая части должны совпасть.

\section*{2. Делитель и кратное: главное различие}
\sectionline

Перед НОД и НОК полезно чётко различать два слова.

\begin{itemize}[leftmargin=7mm,itemsep=3pt,topsep=2pt]
\item \textbf{Делитель числа} --- число, на которое исходное число делится без остатка. Например, \(6\) --- делитель \(24\), потому что \(24:6=4\).
\item \textbf{Кратное числа} --- число, которое делится на исходное число без остатка. Например, \(24\) --- кратное \(6\), потому что \(24:6=4\).
\end{itemize}

\keyidea{Короткая опора}

Если \(d=\nodru(a,b)\), то \(a\) и \(b\) делятся на \(d\).
Если \(m=\nokru(a,b)\), то уже \(m\) делится и на \(a\), и на \(b\).

\section*{3. НОД --- наибольший общий делитель}
\sectionline

\textbf{Определение.}
\(\nodru(a,b)\) --- наибольшее натуральное число, на которое оба числа \(a\) и \(b\) делятся без остатка.

\textbf{Как найти НОД по разложению на простые множители:}

\begin{enumerate}[leftmargin=8mm,itemsep=2pt,topsep=2pt]
\item разложите оба числа на простые множители;
\item найдите множители, которые встречаются в обоих разложениях;
\item возьмите каждый общий множитель столько раз, сколько он встречается в обоих разложениях одновременно;
\item перемножьте выбранные множители.
\end{enumerate}

\textbf{Пример.}
\[
24=2\cdot2\cdot2\cdot3,
\qquad
36=2\cdot2\cdot3\cdot3.
\]
Общие множители: \(2\), ещё один \(2\) и \(3\). Поэтому
\[
\nodru(24,36)=2\cdot2\cdot3=12.
\]
Проверка:
\[
24:12=2,
\qquad
36:12=3.
\]

\keyidea{Смысл НОД в задачах}

НОД особенно часто появляется, когда из нескольких данных количеств нужно \textbf{составить одинаковые группы или части без остатка}, причём требуется максимально возможное число групп или наибольший одинаковый размер части.

\section*{4. НОК --- наименьшее общее кратное}
\sectionline

\textbf{Определение.}
\(\nokru(a,b)\) --- наименьшее натуральное число, которое делится и на \(a\), и на \(b\) без остатка.

Для маленьких чисел НОК иногда видно сразу:
\[
\nokru(2,3)=6,
\qquad
\nokru(2,5)=10.
\]

\textbf{Как найти НОК по разложению на простые множители:}

\begin{enumerate}[leftmargin=8mm,itemsep=2pt,topsep=2pt]
\item разложите оба числа на простые множители;
\item перепишите полностью разложение одного числа;
\item из второго разложения добавьте те множители, которых ещё не хватает;
\item перемножьте полученные множители.
\end{enumerate}

\textbf{Пример.}
\[
12=2\cdot2\cdot3,
\qquad
45=3\cdot3\cdot5.
\]
Берём все множители первого числа: \(2\cdot2\cdot3\).
Из второго разложения одной тройки уже хватает, поэтому добавляем ещё одну \(3\) и множитель \(5\):
\[
\nokru(12,45)
=2\cdot2\cdot3\cdot3\cdot5
=180.
\]
Проверка:
\[
180:12=15,
\qquad
180:45=4.
\]

Ещё один пример:
\[
18=2\cdot3\cdot3,
\qquad
24=2\cdot2\cdot2\cdot3,
\]
поэтому
\[
\nokru(18,24)=2\cdot2\cdot2\cdot3\cdot3=72.
\]

\keyidea{Смысл НОК в задачах}

НОК нужен, когда ищут \textbf{первое общее совпадение} периодических событий или \textbf{наименьшее общее количество}, которое должно делиться на несколько заданных величин.

\clearpage
\section*{5. НОД или НОК: как понять по смыслу}
\sectionline

Слова \textit{«наибольшее»} и \textit{«наименьшее»} полезны как подсказки, но решающим остаётся смысл задачи.

\renewcommand{\arraystretch}{1.35}
\begin{tabularx}{\linewidth}{@{}>{\bfseries}p{0.20\linewidth}YY@{}}
\toprule
 & НОД & НОК \\
\midrule
Что происходит?
& Исходные количества делятся на одинаковые части или группы.
& Ищется общее число или момент, кратный нескольким величинам. \\
Типичный вопрос
& «Сколько максимально одинаковых наборов?»; «Какой наибольший одинаковый размер части?»
& «Когда события впервые совпадут?»; «Какое наименьшее количество подойдёт всем условиям?» \\
Самопроверка
& Ответ должен делить каждое исходное число без остатка.
& Ответ должен делиться на каждое исходное число без остатка. \\
\bottomrule
\end{tabularx}

\subsection*{Пример: одинаковые наборы}

Из \(72\) красных и \(96\) синих бусин нужно составить максимально возможное число одинаковых наборов, используя все бусины.
Число наборов должно делить и \(72\), и \(96\), поэтому нужен НОД:
\[
\nodru(72,96)=24.
\]
Получится \(24\) набора, в каждом по
\[
72:24=3
\]
красные и
\[
96:24=4
\]
синие бусины.

\subsection*{Пример: совпадение периодов}

Моторное масло меняют каждые \(15\) тыс. км, приводной ремень --- каждые \(60\) тыс. км.
Первый пробег после начального момента, при котором обе замены совпадут:
\[
\nokru(15,60)=60\text{ тыс. км}.
\]
Это четвёртая замена масла, потому что \(60:15=4\), и первая плановая замена ремня, потому что \(60:60=1\).

\subsection*{Пример: одинаковая минимальная сумма}

Дима покупает пирожки по \(45\) руб. за штуку, а Витя --- по \(60\) руб. за штуку.
Нужно найти наименьшую одинаковую сумму, которую каждый сможет потратить полностью на целое число пирожков.
Такая сумма должна делиться и на \(45\), и на \(60\):
\[
\nokru(45,60)=180\text{ руб.}
\]
Действительно,
\[
180:45=4,
\qquad
180:60=3.
\]
Значит, Дима купит \(4\) пирожка, а Витя --- \(3\).

\section*{6. Типичные ошибки}
\sectionline

\begin{itemize}[leftmargin=7mm,itemsep=3pt,topsep=2pt]
\item \textbf{Перепутать направление деления.} НОД делит исходные числа; НОК само делится на исходные числа.
\item \textbf{Для НОД взять лишний множитель.} В НОД входят только множители, общие для всех рассматриваемых чисел.
\item \textbf{Для НОК забыть недостающий множитель.} После записи одного разложения нужно добавить из другого всё, чего ещё не хватает.
\item \textbf{Выбрать метод только по одному слову.} Сначала определите смысл: разделение на одинаковые группы или первое общее совпадение.
\item \textbf{Не проверить ответ.} После вычисления НОД или НОК обязательно выполните короткую проверку делением.
\end{itemize}

% ============================================================
% ОТДЕЛЬНАЯ СТРАНИЦА С БЛОК-СХЕМОЙ
% ============================================================

\clearpage
\thispagestyle{empty}

\begin{center}
{\Large\bfseries\color{darkaccent}
Алгоритм: как выбрать НОД или НОК в текстовой задаче\par}
\vspace{11mm}

\begin{tikzpicture}[x=1mm,y=1mm]

\node[flowstart] (start) at (0,0) {Прочитайте условие};

\node[flowaction,text width=50mm] (meaning) at (0,-22)
{Сформулируйте своими словами, что требуется найти};

\node[flowdecision,text width=36mm] (split) at (0,-55)
{Из данных количеств составляют одинаковые группы или части без остатка?};

\node[flowresult,text width=36mm] (gcd) at (-44,-94)
{Используйте НОД};

\node[flowdecision,text width=36mm] (coincide) at (0,-112)
{Ищется первое совпадение или наименьшее общее кратное?};

\node[flowaction,text width=40mm] (gcdcheck) at (-44,-132)
{Проверьте: ответ делит каждое исходное число};

\node[flowresult,text width=36mm] (lcm) at (44,-151)
{Используйте НОК};

\node[flowaction,text width=41mm] (reread) at (-44,-171)
{Перечитайте вопрос и уточните, что именно должно быть общим};

\node[flowaction,text width=40mm] (lcmcheck) at (44,-185)
{Проверьте: ответ делится на каждое исходное число};

\node[flowstart] (finish) at (0,-222)
{Ответ обоснован};

\draw[flowarrow] (start) -- (meaning);
\draw[flowarrow] (meaning) -- (split);

\draw[flowarrow] (split) -- node[above left,font=\small]{да} (gcd);
\draw[flowarrow] (split) -- node[right,font=\small]{нет} (coincide);

\draw[flowarrow] (gcd) -- (gcdcheck);
\draw[flowarrow] (coincide) -- node[above right,font=\small]{да} (lcm);
\draw[flowarrow] (coincide) -- node[above left,font=\small]{нет} (reread);
\draw[flowarrow] (lcm) -- (lcmcheck);

\draw[flowarrow] (gcdcheck.south) |- (finish.west);
\draw[flowarrow] (lcmcheck.south) |- (finish.east);
\draw[flowarrow] (reread.west) -- ++(-10,0) |- (meaning.west);

\end{tikzpicture}
\end{center}

\clearpage

% ============================================================
% ТРЕНИРОВОЧНЫЙ БЛОК
% ============================================================

\section*{Тренировочный блок · Путь А}
\sectionline

\textbf{1. Определите, что следует использовать: НОД или НОК.}

\begin{enumerate}[label=\alph*),leftmargin=8mm,itemsep=2pt,topsep=2pt]
\item Из двух видов конфет составляют максимально возможное число одинаковых наборов без остатка.
\item Два сигнала повторяются через разные промежутки времени. Нужно найти момент их первого нового совпадения.
\item Нужно найти наименьшую одинаковую сумму, которая делится на две заданные цены без остатка.
\end{enumerate}

\vspace{2mm}

\textbf{2.} Разложите числа на простые множители и найдите:
\[
\nodru(30,45),
\qquad
\nokru(8,12).
\]

\vspace{2mm}

\textbf{3.} Есть \(72\) красные и \(96\) синих бусин. Нужно составить максимально возможное число одинаковых наборов, используя все бусины. Сколько наборов получится? Сколько бусин каждого цвета будет в одном наборе?

\vspace{2mm}

\textbf{4.} Два сигнала прозвучали одновременно. Первый повторяется каждые \(18\) секунд, второй --- каждые \(24\) секунды. Через сколько секунд они впервые снова прозвучат одновременно?

\vspace{2mm}

\textbf{5.} Дима покупает пирожки по \(45\) руб. за штуку, а Витя --- по \(60\) руб. за штуку. Какую наименьшую одинаковую сумму нужно иметь каждому, чтобы оба могли потратить её полностью на целое число пирожков? Сколько пирожков купит каждый? Укажите, почему здесь нужен НОД или НОК.

\vfill

{\color{softgray}\small
Для самопроверки выбора метода задайте себе один вопрос:
\textit{я делю исходные количества на одинаковые группы или ищу первое общее совпадение / наименьшее общее кратное?}}

% ============================================================
% ОТВЕТЫ
% ============================================================

\clearpage

\section*{Ответы}
\sectionline

\renewcommand{\arraystretch}{1.4}
\begin{tabularx}{\linewidth}{@{}p{0.07\linewidth}X@{}}
\toprule
\textbf{№} & \textbf{Ответ} \\
\midrule
1
& а) НОД; б) НОК; в) НОК. \\
2
& \(30=2\cdot3\cdot5\), \(45=3\cdot3\cdot5\), поэтому \(\nodru(30,45)=15\). \(8=2\cdot2\cdot2\), \(12=2\cdot2\cdot3\), поэтому \(\nokru(8,12)=24\). \\
3
& \(24\) набора; в каждом по \(3\) красные и \(4\) синие бусины. \\
4
& \(72\) секунды. \\
5
& \(180\) руб.; Дима купит \(4\) пирожка, Витя --- \(3\); используется НОК, потому что искомая сумма должна быть наименьшим общим кратным чисел \(45\) и \(60\). \\
\bottomrule
\end{tabularx}

\vspace{8mm}

\keyidea{Финальная самопроверка}

\begin{itemize}[leftmargin=7mm,itemsep=3pt,topsep=2pt]
\item Для НОД убедитесь, что ответ делит каждое исходное число.
\item Для НОК убедитесь, что ответ делится на каждое исходное число.
\item В текстовой задаче отдельно проверьте, отвечает ли найденное число именно на поставленный вопрос.
\end{itemize}

\end{document}